\subsection{A local properties problem for difference sets}

We study an arithmetic version of such local properties problems, suggested by Zeev Dvir (see \cite{PS19}). We consider sets of real numbers for which every small piece is `arithmetically unstructured,' and we wish to understand how arithmetically unstructured this local property requires the entire set to be. One measure of the extent to which a set is arithmetically structured is the size of its difference set --- roughly, sets with smaller difference set relative to their size possess more arithmetic structure. So we consider sets for which every small subset has a reasonably large difference set, and attempt to understand how large the difference set of the entire set must be. 

We now make this precise. For a set $A \subseteq \RR$, we define the \emph{difference set} of $A$ to be \[A - A = \{\abs{a - b} \mid a, b \in A, \, a \neq b\}.\] Note that we include only positive differences in $A - A$; this definition is somewhat nonstandard but more natural for the problem we study. (When discussing the differences present in a set, we refer only to positive differences, unless otherwise stated.)

For integers $n \geq k \geq 1$ and $\ell \geq 0$, we define $g(n, k, \ell)$ to be the minimum value of $\abs{A - A}$ over all $n$-element sets $A \subseteq \RR$ with the property that every $k$-element subset $A' \subseteq A$ satisfies $\abs{A' - A'} \geq \ell$; we refer to this property as the \emph{$(k, \ell)$-local property}. (Note that the same definition can be used to define $g(n, k, \ell)$ even when $\ell$ is not a nonnegative integer; then $g(n, k, \ell) = g(n, k, \ceil{\ell})$. This makes certain results slightly more convenient to state.) We think of $k$ and $\ell$ as fixed and $n$ as large, and we study the asymptotic behavior of $g(n, k, \ell)$ as $n$ tends to infinity. (In particular, in all asymptotic notation in this paper, $n$ tends to infinity while all other quantities are fixed; the implied constants may depend on $k$, $\ell$, and any other stated parameters.)

As with the local properties problem for graph colorings, we can attempt to understand $g(n, k, \ell)$ either by fixing $\ell$, or by considering various thresholds. Note that for all $n$-element sets $A \subseteq \RR$ we have $n - 1 \leq \abs{A - A} \leq \binom{n}{2}$ (and $\abs{A - A} = n - 1$ if and only if $A$ is an arithmetic progression), so for all $\ell \leq \binom{k}{2}$ we have \[n - 1 \leq g(n, k, \ell) \leq \binom{n}{2}.\] (If $\ell > \binom{k}{2}$, then no set $A$ can satisfy the $(k, \ell)$-local property, so we say that $g(n, k, \ell) = +\infty$.) Also note that $g(n, k, \ell)$ is weakly increasing in $\ell$, as increasing $\ell$ only makes the $(k, \ell)$-local property stronger. So it is natural to consider, for fixed $k$, how long $g(n, k, \ell)$ stays linear and when it becomes quadratic as we increase $\ell$. To formalize this, as defined in \cite{Li22}, the \emph{superlinear threshold} is defined to be the largest integer $\ell$ (as a function of $k$) such that $g(n, k, \ell) = O(n)$, and the \emph{quadratic threshold} is defined to be the smallest integer $\ell$ (as a function of $k$) such that $g(n, k, \ell) = \Omega(n^2)$. (Note that $g(n, k, k - 1) = n - 1$, so at the superlinear threshold we actually have $g(n, k, \ell) = \Theta(n)$; similarly, if $k \geq 4$ then $g(n, k, \binom{k}{2}) = \binom{n}{2}$, so at the quadratic threshold we actually have $g(n, k, \ell) = \Theta(n^2)$.)  
